\documentclass{article}
\usepackage{graphicx} % Required for inserting images
\usepackage{amsmath}

\begin{document}
\subsection{Hierakisk Klyngedannelse}
Klyngedannelse, også kaldet clustering, er en metode inden for matematik og statistik, hvis formål er  at organisere og finde mønstre i data som herefter kan analyseres. Ideen med klyngedannelse, er at opdele et antal observationer i grupper, også kaldet klynger. Den generelle ide er at opdele observationerne i klynger, baseret på hvor meget de "ligner" hindanen. For alle typer klynge-metoder er grupperingen baseret på en valgt afstandsmåling.  Euklidisk-afstand, Manhattan-afstand og korrelationsbaserede afstande kan alle være fornuftige at bruge afhænigt af hvilket type problem det er man vil løse. Dette valg har stor betydning og vil påvirke strukturen, størrelsen af klyngerne, samt beregningskompleksiteten af algoritmen. I alle typer klyngedannelsesmetoder kræves det at tage tre beslutninger:

\begin{itemize}
    \item Valg af afstandsfunktion, $\textit{d}(x_i, x_j)$ 
    \item Et optimeringskriterium eller en strukturregel
    \item En metode til at danne klynger ud fra disse
\end{itemize}

En af de mest udbredte klyngedannelsesalgortimer er "K-means clustering", som er en partitional klyngedannelsesmetode. Her fastlægges et samlet antal klynger (K), hvor formålet er at minimere den samlede variation inden for klyngerne. Denne metode er specielt brugbar når man på forhånd ved hvor mange klynger der skal dannes. Denne metode bliver dermed ofte brugt i optimeringsproblemer hvor man har en kendt begrænsning som svarer til antallet af klynger. Selvom denne metode er effektiv i mange tilfælde er det ikke nødvendigvis den bedste metode at bruge i porteføljeoptimering. Som udgangspunkt kan det være svært at vide hvor mange klynger der skal dannes ud fra et stort datasæt af aktiver og dermed kan det være mere nyttigt at bruge hierarkisk klyngedannelse, hvor denne begrænsing ikke er gældene. Da HRP er bygget op omkring Hierakisk klyngedannelse er det udelukkende denne type vi har fokus på.

Hierarkisk klyngedannelsesaloritmer er enten divisiv eller agglomerativ. Begge metoder har hver deres fordele, men i praksis bruges den agglomerative metode mest i akedemiske sammenhæng. Da HRP også anvender agglomerativ klyngedannelse, fokuserer vi udelukkende på denne metode. Den agglomerative metode (bottom up) starter med enkeltpunkter og iterativt forbinder klynger, indtil alle observationer til sidst sammensættes til en stor klynge, sammensat af alle observationer. Dette er også kaldet "Tree-Clustering" da det grafisk ligner et stamtræ. Mere præcist danner det en graf kaldet et dendrogram.

Formelt set, kan hierakisk klyngedannelse formuleres vha. et datasæt, med punkter givet ved:
\begin{equation}
[
X = {x_1, x_2, \dots, x_n}, \quad x_i \in \mathbb{R}^d
]
\end{equation}
Lad \begin{equation}
    (D = (d_{ij})) 
\end{equation}
være en afstandsmatrice, hvor
\begin{equation}
[
d_{ij} = d(x_i, x_j)
]
\end{equation}
er en metrisk afstand.\\

Man starter med (n) klynger:
\begin{equation}
[
C_1 = {x_1}, \dots, C_n = {x_n}
]
\end{equation}
Herefter finder man de to klynger (C_a) og (C_b), der har kortest/mindst afstand:
\begin{equation}
[
(C_a, C_b) = \arg\min_{C_i \neq C_j} d(C_i, C_j)
]
\end{equation}
Man erstater dem med den samlede klynge:
\begin{equation}
[
C_{ab} = C_a \cup C_b
]
\end{equation}
Man opdaterer afstandene og beregner nu afstanden mellem den nye klynge og de øvrige.

Man gentager nu processen indtil alle observationer er samlet i en klynge.

\subsubsection{Klyngeafstandsmål}
\end{document}